\section{Classifier Guidance：分类器引导}

\subsection{从联合分布到条件分布}

条件生成的目标是学习条件分布 $p(x|y)$，其中 $y$ 是条件信息（如类别标签）。关键洞察在于，联合分布和条件分布的 score function 之间存在简单关系：

$$
\nabla_{x_t}\log p(x_t, y) = \nabla_{x_t}\log p(x_t) + \nabla_{x_t}\log p(y|x_t)
$$

这是因为 $p(x_t, y) = p(x_t)\cdot p(y|x_t)$，取对数后对 $x_t$ 求梯度，$\log p(y)$ 项为关于 $x_t$ 的常数而消失。

\begin{importantbox}{Classifier Guidance 核心公式}
在扩散模型的每一步去噪中，修改 score function：
$$
\tilde{\nabla}_{x_t}\log p(x_t, y) = \underbrace{\nabla_{x_t}\log p(x_t)}_{\text{无条件 score（预训练模型）}} + \underbrace{\nabla_{x_t}\log p(y|x_t)}_{\text{分类器梯度引导}}
$$
对应到噪声预测，更新规则为：
$$
\tilde{\epsilon}(x_t, t) = \epsilon_\theta(x_t, t) - \sqrt{1-\bar\alpha_t}\;\nabla_{x_t}\log p_\phi(y|x_t)
$$
其中 $p_\phi(y|x_t)$ 是在带噪数据上训练的分类器。
\end{importantbox}

\subsection{噪声分类器的训练}

这里有一个关键细节：分类器 $p_\phi(y|x_t)$ 不是普通的图像分类器，而是需要能处理\textbf{任意噪声水平}的中间数据点 $x_t$。训练方法如下：

\begin{enumerate}
    \item 对于数据集中的每一对 $(x_0, y)$，随机采样时间步 $t$
    \item 通过前向过程 $q(x_t|x_0)$ 得到带噪图像 $x_t$
    \item 训练分类器以 $x_t$ 和 $t$ 为输入，预测类别 $y$ 的概率
\end{enumerate}

\begin{warningbox}{分类器必须处理带噪输入}
不能直接使用在干净图像上训练的 ImageNet 分类器！分类器必须在各种噪声水平的 $x_t$ 上训练，否则梯度信号在高噪声步骤时完全不可靠。这是 Classifier Guidance 的一个额外训练成本。
\end{warningbox}

\subsection{引导强度的控制}

通过引入权重 $w$ 来控制条件的强度：

$$
\tilde{\epsilon}(x_t, t) = \epsilon_\theta(x_t, t) - w\cdot\sqrt{1-\bar\alpha_t}\;\nabla_{x_t}\log p_\phi(y|x_t)
$$

\begin{itemize}
    \item $w = 0$：纯无条件生成
    \item $w = 1$：标准条件生成（对应联合分布的 score）
    \item $w > 1$：增强条件效果——输出更忠于条件，但多样性降低
\end{itemize}

\begin{knowledgebox}{引导强度的直觉理解}
从轨迹的角度看，无条件扩散模型的去噪轨迹从 $x_T$ 到 $x_0$ 遵循无条件 score field。增加分类器梯度项相当于在每一步\textbf{偏转}轨迹，使其朝向满足条件 $y$ 的数据流形区域移动。增大 $w$ 使偏转更剧烈，输出更"典型"但更缺乏多样性。
\end{knowledgebox}

\subsection{Classifier Guidance 的优势与不足}

\begin{table}[H]
\centering
\begin{tabular}{p{6cm}p{6cm}}
\toprule
\textbf{优势} & \textbf{不足} \\
\midrule
可复用预训练的无条件模型 & 需要额外训练噪声分类器 \\
灵活控制引导强度 & 推理时需运行两个网络 \\
不需要重训扩散模型 & 分类器梯度可能不稳定 \\
数学推导清晰 & 仅限于有分类器可用的条件类型 \\
\bottomrule
\end{tabular}
\caption{Classifier Guidance 优缺点对比}
\end{table}

这些不足直接促成了 Classifier-Free Guidance 的提出。

\subsection{本章小结}

Classifier Guidance 的核心思想是：利用贝叶斯分解，将条件生成分解为无条件 score 加上分类器梯度。这一思想揭示了扩散模型相较于 GAN 和 VAE 的独特优势——可以在\textbf{不重训}生成模型的情况下，通过修改 score function 实现条件引导。这种 score manipulation 的范式将贯穿整个课程的后续内容。

\newpage
